% answers_short.tex —— 2005 年·普通高等学校招生考试·山东卷·文科 · 简短答案速查
% 由 tools/make_short_answers.py 自 answers.tex 的 \gksolhead/\gkanswer 提取，
% 并与 2005/2005_数学真题_山东文科_答案要点.md 逐题比对后生成，勿手改（改源 MD 后重跑）。

\begin{gknotice}
\textbf{使用说明}：本节只给最终答案便于快速核对，完整推导见解析版。答案为编写者独立推导，\textbf{非官方标准答案}。
\end{gknotice}

\gkpart{选择题}
\begin{center}\begin{tabular}{ccccccccccccc}
\toprule
题号 & 1 & 2 & 3 & 4 & 5 & 6 & 7 & 8 & 9 & 10 & 11 & 12 \\
\midrule
答案 & C & C & B & B & D & C & C & A & D & D & A & B \\
\bottomrule
\end{tabular}\end{center}

\gkpart{填空题}
\begin{itemize}[leftmargin=2.6em]
  \item[\textbf{13.}] $50$。
  \item[\textbf{14.}] $\sqrt{2}$。
  \item[\textbf{15.}] $(2, 3)$。
  \item[\textbf{16.}] ③④。
\end{itemize}

\gkpart{解答题}
\begin{itemize}[leftmargin=2.6em]
  \item[\textbf{17.}] $-\dfrac{4}{5}$。
  \item[\textbf{18.}] (1) 袋中原有白球 $3$ 个；(2) 取球 $2$ 次终止的概率为 $\dfrac{2}{7}$；(3) 甲取到白球的概率为 $\dfrac{22}{35}$
  \item[\textbf{19.}] (1) $n = 3m + 6$；(2) 函数 $f(x)$ 的单调递增区间为 $\left(1 + \dfrac{2}{m}, 1\right)$；单调递减区间为 $\left(-\infty, 1 + \dfrac{2}{m}\right)$ 和 $(1, +\infty)$
  \item[\textbf{20.}] (1) 异面直线 $AE$ 与 $BF$ 所成的角为 $\arccos\dfrac{\sqrt{2}}{4}$；(2) 平面 $BDF$ 与平面 $AA_1B$ 所成的二面角（锐角）的大小为 $\arccos\dfrac{\sqrt{15}}{5}$；(3) 点 $A$ 到平面 $BDF$ 的距离为 $\dfrac{2\sqrt{5}}{5}$
  \item[\textbf{21.}] (1) 见解析版；(2) $f'(1) = 3(n-1)2^{n+1} + 6 - \dfrac{n(n+1)}{2}$
  \item[\textbf{22.}] (1) 动圆圆心 $C$ 的轨迹方程为 $y^2 = 2px$；(2) 见解析版
\end{itemize}
